\begin{figure}[htbp]
\centering
\begin{tikzpicture}[font=\small,>=Stealth]

\filldraw[fill=blue!8,draw=blue] (-1,-.6)--(2,-.6)--(3,.4)--(0,.4)--cycle;
\fill (1,0) circle (2pt);
\draw[->,very thick,blue] (1,0)--(1,2.3) node[above] {$n$};
\draw[->,very thick,red!70!black] (1,0)--(3.2,1.6) node[right] {$F$};
\draw[dashed] (3.2,1.6)--(1,1.6);
\draw[->,thick,red!70!black] (.88,0)--(.88,1.6) node[left] {$(F\cdot n)n$};
\node[align=left,anchor=west] at (4,1) {$\displaystyle n=\dfrac{\sigma_u\times\sigma_v}{\|\sigma_u\times\sigma_v\|}$\\[8pt]
$F\cdot n\,dS$\\[3pt]$\quad=F\cdot(\sigma_u\times\sigma_v)\,du\,dv$};
\node[align=center] at (4,-1.3) {Swap $u,v$: the normal and flux reverse;\\the parallelogram area is unchanged.};

\end{tikzpicture}
\caption{Flux measures the normal component of the field. Tangential motion contributes no flux through the patch. The sketch uses a positively directed normal component.}
\label{fig:flux-normal}
\end{figure}
